% TOP_PROBLEM: {"problemId": "problem.anisotropic-calderon-conjecture-for-smooth-riemannian-manifolds", "canonicalTitle": "Smooth anisotropic Calderón uniqueness from full boundary data", "releaseRank": 176, "primaryDomain": "analysis_pde", "primaryDomainLabel": "Analysis & PDE", "displayStatus": "open_with_solved_subcases", "releaseStatus": "open_with_solved_subcases"}
% !TeX program = lualatex
% Definition and sources only; this file is not an LLM solution attempt.
\documentclass[11pt]{article}
\usepackage[margin=25mm]{geometry}
\usepackage{fontspec}
\setmainfont{Latin Modern Roman}
\usepackage{amsmath,amssymb,mathtools}
\usepackage{unicode-math}
\setmathfont{Latin Modern Math}
\usepackage{xurl,hyperref}
\hypersetup{colorlinks=true,urlcolor=blue}
\setcounter{secnumdepth}{0}
\setlength{\parindent}{0pt}
\setlength{\parskip}{5pt}
\setlength{\emergencystretch}{3em}
\begin{document}
\section*{\texorpdfstring{Smooth anisotropic Calderón uniqueness from full boundary data}{Smooth anisotropic Calderón uniqueness from full boundary data}}\label{176}
\subsection{Definitions and mathematical statement}
For a smooth metric $g$ on a compact connected manifold $M$ with boundary, solve the Dirichlet problem $\Delta_g u=0$, $u|_{\partial M}=f$, and define the Dirichlet-to-Neumann operator
$\Lambda_g f=\partial_{\nu_g}u|_{\partial M}$, with outward metric-unit normal $\nu_g$. In dimension $n\ge3$, the target is
\[
 \Lambda_{g_1}=\Lambda_{g_2}\quad\Longrightarrow\quad
 g_2=\Phi^*g_1\text{ for some }\Phi|_{\partial M}=\mathrm{id}.
\]
Both metrics live on the same smooth manifold and the operators are compared on the entire boundary using its identity marking. Pullback by a boundary-fixing diffeomorphism is the unavoidable gauge ambiguity. No real-analyticity, product structure, or partial-boundary restriction is included.
\par\smallskip\textit{Formulation and concept references:} \href{https://arxiv.org/pdf/2608.15813v1}{[S1]}.

\subsection{Short English statement}
Do all boundary voltage-to-flux measurements determine a smooth interior metric, apart from a change of coordinates that fixes the boundary?
\subsection{Sources}
\begin{itemize}
\item[S1]Quasianalyticity and geometric rigidity in anisotropic Calderon's problem; arXiv2608.15813v1 August16,2026.
\url{https://arxiv.org/pdf/2608.15813v1}.
\textit{Formulation source.}
\end{itemize}
\end{document}
